\documentclass{article}
\usepackage{hyperref}
\usepackage{/globals/anki}
\ankiprefix{mengen-umformungen}
\ankitopic{Mengen Umformungen}
\ankishow
\begin{document}
\section{Äquivalenzumformungen für Mengen} 
Die Rechenregeln der Mengen sind eigentlich den \hyperref[sec:logik_umformungen]{Regeln der Äquivalenzumformungen der Aussagenlogik} identisch. Dabei ist eine Komplement eine Negation, eine Vereinigung ein Disjunktion, ein Durchschnitt eine Konjunktion, eine Grundmenge eine 1 und eine leere Menge eine 0.
\ankinote{vereinigung}{Vereinigung}{Disjunktion}
\ankinote{durchschnitt}{Durchschnitt}{Konjunktion}

Desweiteren gilt
\[
 A \backslash B \quad = \quad A \cap \overline{B}
\]
\ankinote{a-minus-b}{\(A \backslash B \quad\)}{\(\quad A \cap \overline{B}\)}

Mit diesen Regeln können Formeln über Mengen vereinfacht werden.
\end{document}